% GENERATED FILE. Edit canonical lessons, then run npm run generate:books.
% canonical-source-hash: fnv1a64:2caf497065aa29fc
% canonical-lessons: GU-C342-vivad, GU-C342-takrar, GU-C342-matbhed, GU-C342-khandan, GU-C342-tokvu

\chapter{A Dispute}
\label{ch:gu-a-dispute}
\hlchaptermodality{342}

\begin{chapteropening}
\textbf{By the end of this chapter:} \emph{I can say what two sides disagreed about.}

Complete the last lesson of chapter 342: I can say what two sides disagreed about.
\end{chapteropening}

\section[\texorpdfstring{\gu{વિવાદ} (vivād)}{વિવાદ (vivād)}]{\gu{વિવાદ} (vivād) — a dispute}

\label{lesson:GU-C342-vivad}



\begin{quote}
Before the new one: say the Gujarati for a context, then the Gujarati for to raise a point.
\end{quote}

\subsection*{You'll want to know: \gu{વિવાદ}}
\textbf{\gu{વિવાદ}} — \emph{vivād} — \textquotedblleft{}a dispute\textquotedblright{}.

\begin{quote}
\gu{મિલકત} \gu{માટે} \gu{બે} \gu{ભાઈઓ} \gu{વચ્ચે} \gu{વિવાદ} \gu{થયો}.
\end{quote}

\begin{quote}
\emph{milkat māṭe be bhāīo vacce vivād thayo.}
\end{quote}

\begin{quote}
A dispute arose between two brothers over property.
\end{quote}

A masculine noun.

\begin{cousinweb}
\textbf{\gu{વિવાદ}} (\emph{vivād}) is a tatsama, borrowed from Sanskrit \emph{vivāda}, \textquotedblleft{}a dispute, a quarrel\textquotedblright{} (Monier-Williams).
\end{cousinweb}

\begin{grammarlens}[title={What you've built}]
One more word for giving reasons and opinions.
\end{grammarlens}

\subsection*{Your turn}
\begin{itemize}
\raggedright
  \item \emph{Say it:} \emph{vivād}
  \item \emph{Say it:} \emph{vivād}, once more
  \item \emph{Say it:} \emph{milkat māṭe be bhāīo vacce vivād thayo.} — A dispute arose between two brothers over property.
  \item \emph{Recall:} say the Gujarati for a context, then the Gujarati for to raise a point, then say \emph{vivād} again
\end{itemize}

\begin{tcolorbox}[breakable,colback=teal!4,colframe=teal!35!black,title={Before you move on}]
What does \gu{વિવાદ} mean? (\textquotedblleft{}A dispute\textquotedblright{}.) Say it once more.
\end{tcolorbox}
\section[\texorpdfstring{\gu{તકરાર} (takrār)}{તકરાર (takrār)}]{\gu{તકરાર} (takrār) — a squabble}

\label{lesson:GU-C342-takrar}



\begin{quote}
Before the new one: say the Gujarati for to raise a point, then the Gujarati for a dispute.
\end{quote}

\subsection*{You'll want to know: \gu{તકરાર}}
\textbf{\gu{તકરાર}} — \emph{takrār} — \textquotedblleft{}a squabble\textquotedblright{}.

\begin{quote}
\gu{બાળકો} \gu{રમકડા} \gu{માટે} \gu{તકરાર} \gu{કરતાં} \gu{હતાં}.
\end{quote}

\begin{quote}
\emph{bāḷko ramakḍā māṭe takrār kartā̃ hatā̃.}
\end{quote}

\begin{quote}
The children were squabbling over a toy.
\end{quote}

A feminine noun.

\begin{cousinweb}
\textbf{\gu{તકરાર}} (\emph{takrār}) is usually traced to Arabic \emph{takrār}, \textquotedblleft{}repetition\textquotedblright{}, through Persian; Gujarati uses it for a petty quarrel, a sense the donor word does not have.
\end{cousinweb}

\begin{grammarlens}[title={What you've built}]
One more word for giving reasons and opinions.
\end{grammarlens}

\subsection*{Your turn}
\begin{itemize}
\raggedright
  \item \emph{Say it:} \emph{takrār}
  \item \emph{Say it:} \emph{takrār}, once more
  \item \emph{Say it:} \emph{bāḷko ramakḍā māṭe takrār kartā̃ hatā̃.} — The children were squabbling over a toy.
  \item \emph{Recall:} say the Gujarati for to raise a point, then the Gujarati for a dispute, then say \emph{takrār} again
\end{itemize}

\begin{tcolorbox}[breakable,colback=teal!4,colframe=teal!35!black,title={Before you move on}]
What does \gu{તકરાર} mean? (\textquotedblleft{}A squabble\textquotedblright{}.) Say it once more.
\end{tcolorbox}
\section[\texorpdfstring{\gu{મતભેદ} (matbhed)}{મતભેદ (matbhed)}]{\gu{મતભેદ} (matbhed) — a disagreement}

\label{lesson:GU-C342-matbhed}



\begin{quote}
Before the new one: say the Gujarati for a dispute, then the Gujarati for a squabble.
\end{quote}

\subsection*{You'll want to know: \gu{મતભેદ}}
\textbf{\gu{મતભેદ}} — \emph{matbhed} — \textquotedblleft{}a disagreement\textquotedblright{}.

\begin{quote}
\gu{આ} \gu{યોજના} \gu{પર} \gu{અરુણ} \gu{અને} \gu{મીરા} \gu{વચ્ચે} \gu{મતભેદ} \gu{છે}.
\end{quote}

\begin{quote}
\emph{ā yojnā par Aruṇ ane Mīrā vacce matbhed che.}
\end{quote}

\begin{quote}
Arun and Mira disagree about this plan.
\end{quote}

A masculine noun.

\begin{cousinweb}
\textbf{\gu{મતભેદ}} (\emph{matbhed}) is built from \textbf{\gu{મત}} (\emph{mat}), \textquotedblleft{}an opinion\textquotedblright{}, and Sanskrit \emph{bheda}, \textquotedblleft{}a split, a difference\textquotedblright{} (Monier-Williams): a split of opinion.
\end{cousinweb}

\begin{grammarlens}[title={What you've built}]
One more word for giving reasons and opinions.
\end{grammarlens}

\subsection*{Your turn}
\begin{itemize}
\raggedright
  \item \emph{Say it:} \emph{matbhed}
  \item \emph{Say it:} \emph{matbhed}, once more
  \item \emph{Say it:} \emph{ā yojnā par Aruṇ ane Mīrā vacce matbhed che.} — Arun and Mira disagree about this plan.
  \item \emph{Recall:} say the Gujarati for a dispute, then the Gujarati for a squabble, then say \emph{matbhed} again
\end{itemize}

\begin{tcolorbox}[breakable,colback=teal!4,colframe=teal!35!black,title={Before you move on}]
What does \gu{મતભેદ} mean? (\textquotedblleft{}A disagreement\textquotedblright{}.) Say it once more.
\end{tcolorbox}
\section[\texorpdfstring{\gu{ખંડન} (khanḍan)}{ખંડન (khanḍan)}]{\gu{ખંડન} (khanḍan) — a refutation}

\label{lesson:GU-C342-khandan}



\begin{quote}
Before the new one: say the Gujarati for a squabble, then the Gujarati for a disagreement.
\end{quote}

\subsection*{You'll want to know: \gu{ખંડન}}
\textbf{\gu{ખંડન}} — \emph{khanḍan} — \textquotedblleft{}a refutation\textquotedblright{}.

\begin{quote}
\gu{અરુણે} \gu{મીરાની} \gu{દલીલનું} \gu{ખંડન} \gu{કર્યું}.
\end{quote}

\begin{quote}
\emph{Aruṇe Mīrānī dalīlnũ khanḍan karyũ.}
\end{quote}

\begin{quote}
Arun refuted Mira's argument.
\end{quote}

A neuter noun.

\begin{cousinweb}
\textbf{\gu{ખંડન}} (\emph{khanḍan}) is a tatsama, borrowed from Sanskrit \emph{khaṇḍana}, \textquotedblleft{}breaking, cutting to pieces\textquotedblright{} (Monier-Williams): an argument broken down.
\end{cousinweb}

\begin{grammarlens}[title={What you've built}]
One more word for giving reasons and opinions.
\end{grammarlens}

\subsection*{Your turn}
\begin{itemize}
\raggedright
  \item \emph{Say it:} \emph{khanḍan}
  \item \emph{Say it:} \emph{khanḍan}, once more
  \item \emph{Say it:} \emph{Aruṇe Mīrānī dalīlnũ khanḍan karyũ.} — Arun refuted Mira's argument.
  \item \emph{Recall:} say the Gujarati for a squabble, then the Gujarati for a disagreement, then say \emph{khanḍan} again
\end{itemize}

\begin{tcolorbox}[breakable,colback=teal!4,colframe=teal!35!black,title={Before you move on}]
What does \gu{ખંડન} mean? (\textquotedblleft{}A refutation\textquotedblright{}.) Say it once more.
\end{tcolorbox}
\section[\texorpdfstring{\gu{ટોકવું} (ṭokvũ)}{ટોકવું (ṭokvũ)}]{\gu{ટોકવું} (ṭokvũ) — to cut someone short}

\label{lesson:GU-C342-tokvu}



\begin{quote}
Before the new one: say the Gujarati for a disagreement, then the Gujarati for a refutation.
\end{quote}

\subsection*{You'll want to know: \gu{ટોકવું}}
\textbf{\gu{ટોકવું}} — \emph{ṭokvũ} — \textquotedblleft{}to cut someone short\textquotedblright{}.

\begin{quote}
\gu{અરુણ} \gu{દરેક} \gu{વખતે} \gu{મીરાને} \gu{ટોકે} \gu{છે}.
\end{quote}

\begin{quote}
\emph{Aruṇ darek vakhte Mīrāne ṭoke che.}
\end{quote}

\begin{quote}
Arun cuts Mira short every time.
\end{quote}

Past, with \gu{તેણે} (he or she …): \gu{તેણે} \gu{છોકરાને} \gu{ટોક્યો}, \emph{teṇe chokrāne ṭokyo}.

\begin{cousinweb}
\textbf{\gu{ટોકવું}} (\emph{ṭokvũ}) is inherited from \emph{ṭōkk}, \textquotedblleft{}strike, bite\textquotedblright{}, a form CDIAL reconstructs; CDIAL glosses the Gujarati verb \textquotedblleft{}to scold\textquotedblright{} (CDIAL 5476), and today it is said of cutting a speaker short.
\end{cousinweb}

\begin{grammarlens}[title={What you've built}]
One more word for giving reasons and opinions.
\end{grammarlens}

\subsection*{Your turn}
\begin{itemize}
\raggedright
  \item \emph{Say it:} \emph{ṭokvũ}
  \item \emph{Say it:} \emph{ṭokvũ}, once more
  \item \emph{Say it:} \emph{Aruṇ darek vakhte Mīrāne ṭoke che.} — Arun cuts Mira short every time.
  \item \emph{Recall:} say the Gujarati for a disagreement, then the Gujarati for a refutation, then say \emph{ṭokvũ} again
\end{itemize}

\begin{tcolorbox}[breakable,colback=teal!4,colframe=teal!35!black,title={Before you move on}]
What does \gu{ટોકવું} mean? (\textquotedblleft{}To cut someone short\textquotedblright{}.) Say it once more.
\end{tcolorbox}
